


\clearpage
%\setcounter{page}{21}

\section {Arhitectura sistemului }
Arhitectura sistemului este una distribuită integrând concepte de IoT.  Aceasta structură are ca scop asigurarea unui control determinist și în timp real al unui braț robotic fizic prin intermediul unui mediu de realitate virtuală.

 Sistemul a fost realizat din 4 niveluri funcționale principale, care comunică pentru a transforma controlul provenit de la utilizator în mișcarea mecanică a brațului real.

\subsection{Nivelul de simulare și interacțiune}
Acest nivel rulează folosind casca de realitate virtual cu aplicația efectivă rulând pe calculatorul principal, fiind dezvoltată folosind motorul grafic Unity 3D. Rolul său este de a prelua comenzile provenite de la utilizator, de a randa mediul 3D și de a procesa matematica și controlul asociat mișcării. Componentele de bază ale acestui nivel sunt:


\begin{itemize}
	\item Subsistemul VR și Gestionarea comenzilor: Implementat folosind pachetul "XR Interacțion Toolkit". Aceasta interfațează direct cu driverele căștii pentru a citi poziția spațială pentru cap și manete. Scriptul "ToggleXRControls" gestionează stările de control, în timp ce maparea butoanelor fizice este realizată prin asocieri în profilele de intrări prin "XRI Default Input Actions".

	\item Modelul cinematic și "Digital Twin": Reprezentarea virtuală a robotului se bazează pe importul fișierului de descriere URDF. Acesta definește ierarhia rigidă a componentelor și a tipului articulațiilor. Geometria vizuală este asigurată prin randarea 3D a fișierelor de tip STL.

	\item Controlul direct :Reprezintă o modalitate de  mișcare a articulațiilor în care comanda care provenea de la utilizator se aplica direct pe articulațiile principale ale brațului. 

	\item Modulul de cinematica inversă: Mișcarea robotului virtual provine de la script-ul de cinematică inversă și cel de control al efectorului final. Atunci când utilizatorul folosește maneta, scriptul de control a efectorului final actualizează poziția unui obiect din scenă. Script-ul de cinematică inversă preia coordonatele carteziene X, Y, Z ale țintei și calculează unghiurile de rotație necesare pentru fiecare articulație
	\item Elemente ajutătoare de control: Implementat prin citirea de JSON al fișierului care conține lista de controale și acțiunile pe care utilizatorul le poate face asupra robotului. Conține și indicatoare pentru a marca dacă acțiunea respectivă a fost realizată.
\end{itemize}

\subsection{Nivelul de comunicație}
Acest nivel asigură legătura bidirecțională wireless între Unity și robotul fizic. Componentele de bază ale acestui nivel:
\begin{itemize}
	\item Canalul MQTT: Pentru transmiterea comenzilor, a fost implementat protocolul MQTT ce folosește ca model o arhitectură de tip "Publish/Subscribe"
	\item Tunelul prin "ngrok" prin care se poate realiza o conexiune din orice rețea între Unity și ESP
	\item Fluxul de comandă dinte Unity și ESP: Este folosit un script care preia unghiurile de la brațul virtual, îl serializează într-un pachet de date și le publică pe brokerul MQTT sub un topic dedicat
	\item Fluxul de Feedback dintre ESP și Unity: Este folosit un script în Unity ce ascultă pe topicuri ce conțin mesaje provenite de la robotul real legate de starea curentă și confirmări de execuție ale mișcării
\end{itemize}

\subsection{Nivelul de procesare al robotului fizic}
Sistemul de procesare este amplasat pe șasiul robotului și care folosește un microcontroler ESP32, programat folosind limbajul de programare Micropython. Componentele de bază ale acestui nivel:
\begin{itemize}
	\item Inițializarea și gestiunea rețelei: La inițializarea sistemului, este apelat un script ce se ocupă de configurația rețelei microcontrolerului prin conectarea acestuia la un punct de acces local și pentru  atribuirea unei adrese IP statică. Pe urmă utilizând o implementare pe baza protocolului MQTT se stabilește o sesiune TCP cu brokerul local.
	\item Gestionatorul datelor și parsarea datelor: Scriptul "main" rulează bucla de execuție a procesorului. Acesta comunică direct cu script-ul de gestionare MQTT al microcontrolerului. De fiecare dată când brokerul MQTT livrează un pachet trimis de Unity, script-ul de gestionare MQTT interceptează pachetele, se ocupă de deserializarea acestora și convertește valorile text în variabile care să fie trimise către motoare
	
\end{itemize}


\subsection{Nivelul de execuție hardware}
Acest nivel reprezintă interfața fizică a sistemului, unde semnalele digitale procesate de ESP sunt convertite în semnale mecanice și feedback local. Pentru implementarea acestui nivel producătorul a oferit o serie de scripturi și anume:
\begin{itemize}
	\item Pentru controlul servomotoarelor seriale: Brațul robotic  real folosește servomotoare de magistrală conectate în topologie de tip lanț, printr-un fir de comunicație UART. Scriptul conține logica de comunicare cu controllerul intern al motoarelor. Pentru fiecare mișcare, scriptul generează un cadru binar de date în format hexazecimal care conține: un header de sincronizare, ID-ul unic al motorului, poziția unghiulară a motorului convertită la valori între 0 și 1000, timpul alocat mișcării și un octet de paritate pentru detecția erorilor de transmisie serială.
	\item Controlul motoarelor prin PWM: Pentru mișcări auxiliare, există un script care generează semnale de tip PWM utilizând timer-ele hardware ale ESP-uluu prin ajustarea factorului de umplere.
	\item Sistemul de prindere prin vacuum: Efectorul final al brațului folosește  o pompă de vid pentru manipularea obiectelor. Script-ul robotului ce se ocupă de pompă mapează semnalul boolean primit din Unity. La primirea valorii logice de "1", scriptul activează pompa.
	\item Sistemul de feedback și diagnostic local: Producătorul a oferit posibilitatea de implementa un asemenea sistem direct pe placa ESP prin folosirea unui buzzer și a unui led.
\end{itemize}


\begin{figure}[H]
	\centerline{\includegraphics[width=\columnwidth,keepaspectratio]{ArhitecturaLicenta1_cropped.pdf}}
	\caption[(Arhitectură sistem)]{Arhitectura sistemului 
	}
	\label{fig:top}
\end{figure}
